\chapter*{Preface}
The aim of this guide is to give a reference design flow for the creation of a project using both Vivado and Petalinux SDK. Within this way of programming, the bare metal approach, such as the one normally used for STM micro controllers, gives more space to a OS-oriented method. To all intents and purposes, after the creation of the project, what is left is a tailor-made OS that will manage the board in a similar way to that of the bare metal, but with all the advantages of having a OS.
\bigskip
\\Even though some passages are explained, it is given for granted that the developer has some basic knowledge of the Vivado and Petalinux environments. While a the creation of the Vivado- Petalinux project and further implementations are illustrated, it is left for the reader to find a way to create the environment needed to follow the steps disclosed in the following. To write this document and to create the project the following software/devices have been used:
\begin{itemize}
	\item Kria KR260 Robotic Starter Kit
	\item Vivado 2024.1
	\item Petaliux SDK v2022.2
	\item CenOS stream 9
\end{itemize}
Newest software releases might not be compatible with the procedure(s) listed and/or with the files used to support the processes. It is recommended to check resources availability and compatibility before the installation of Vivado and, paying more attention, of Petalinux SDK.
\bigskip
\\The guide is divided into four chapters. The first chapter gives an overview of the Petalinux installation procedure with the bespoke OS (Operative System). The second chapter deals with the creation of a simple project that will exploit the programmable logic (PL). This method is called bare-metal programming. The third chapter proposes the same project under a new light: the creation of a bootable SD card and the running of the project from within the SD card. The last chapter deals with the creation of a PS application so as to read a register. At the end of each project there is also a debugging session utilizing the J-TAG interface.